% 자동 생성 — scripts/build_tables.py. 손으로 고치지 않는다.
% 출처: docs/decisions/*.md
\begin{table}[htbp]
  \centering
  \caption{결정기록. 설계를 정할 때마다 왜 그렇게 정했는지를 한 편씩 남겼고 근거와 함께 \textbf{버린 대안}도 적었다.}
  \label{tab:decisions}
  \begin{tabular}{llp{9.5cm}}
    \toprule
    \textbf{번호} & \textbf{날짜} & \textbf{제목} \\
    \midrule
    0001 & 09-03 & 주 결과변수를 가격에서 거래량으로 변경 \\
    0002 & 09-03 & 분석 구간을 2025-10-19 이전으로 한정 \\
    0003 & 09-03 & 단지(aptSeq) 고정효과를 기본 사양으로 채택 \\
    0004 & 09-03 & 취소거래 제외 및 신고지연 절단 \\
    0005 & 09-03 & 처치 역전 구간에 de Chaisemartin \& D'Haultfœuille 추정량 사용 \\
    0006 & 09-03 & 모든 결과를 반사실 대비로 명시하여 서술한다 \\
    0007 & 09-07 & 문서를 동결·살아있는·최종으로 나누고, 숫자에 조건을 붙인다 \\
    0008 & 09-07 & 지정 유지 단지를 단지명이 아니라 지번으로 확정한다 \\
    0009 & 09-07 & 통제군 오염 점검 결과와 대응 \\
    0010 & 09-07 & 2020년 지정 범위를 법정동 단위 처치로 확정한다 \\
    0011 & 09-08 & 통제군을 어떻게 고르는가 \\
    0012 & 09-08 & 처치군을 무엇으로 하고, 무엇을 확장으로 미루는가 \\
    0013 & 09-11 & 사건 시점을 효력일에 맞추고, 거래 없는 달을 0으로 채운다 \\
    0014 & 09-11 & 거래량의 사전 추세 위배를 어떻게 다루는가 \\
    0015 & 09-11 & 대출 규제를 어떻게 다루는가 — 무엇이 상쇄되고 무엇이 남는가 \\
    0016 & 09-11 & 위약 검정 — 효과는 있으나 크기는 구간으로 말한다 \\
    0017 & 09-11 & 합성대조군 — 평균 회귀 의심에 답한다 \\
    0018 & 09-11 & 반복매매지수를 어떻게 짓는가 \\
    0019 & 09-18 & 클러스터가 적을 때의 추론 — 관행 p 를 버린다 \\
    0020 & 09-18 & 오염 동을 빼도 직후 효과는 그대로이고, 24개월 효과는 커진다 \\
    0021 & 09-18 & 15억에는 계단이 없다 — 가격 차이는 면적 면제와 겹친다 \\
    0022 & 09-18 & 이질성은 대출이 아니라 허가 면적 기준에서 왔다 \\
    0023 & 09-23 & 2025년 구간은 해제가 아니라 재지정으로 본다 \\
    0024 & 09-23 & 파급은 같은 동 안에만 있었다 — '인근 급등'은 기준점이 만든 말이다 \\
    0025 & 09-23 & 가격 효과는 −4 \textasciitilde{} −7\% 로 모인다 — 다만 합성대조군으로는 확인되지 않는다 \\
    0026 & 09-30 & 사전 추세를 감안하면 직후 거래량만 남는다 \\
    0027 & 09-30 & 음의 가중치는 작다 — 부과는 뚜렷하고 해제는 가릴 수 없다 \\
    0028 & 09-30 & 192개 사양 전부 음수다 — 갈린 것은 통제군이 아니라 날짜 세는 법이었다 \\
    0029 & 10-01 & 폐지 뒤 고가 거래는 돌아왔지만 그 규제 덕으로 돌릴 수 없다 \\
    0030 & 10-01 & 취소거래는 정책 반응이 아니다 — 그리고 0004 의 한계는 생각보다 작다 \\
    0031 & 10-08 & 강남·송파 밖에서도 같다 — B그룹이 A그룹과 다르다고 할 수 없다 \\
    \bottomrule
  \end{tabular}

  \begin{minipage}{\linewidth}\footnotesize
  \textit{주.} 날짜는 \textbf{작성일}이다 --- 뒤에 보강·정정된 기록이 있다. \textbf{이 표는 \texttt{docs/decisions/}의 파일들에서 생성된다.} 손으로 적었을 때는 색인의 날짜 셋이 정작 그 기록과 달랐다.
  \end{minipage}
\end{table}
